% Alur penanganan insiden dari deteksi sampai tindak lanjut, beserta empat peran
% yang dipisahkan selama insiden berlangsung.
\begin{tikzpicture}[
  font=\scriptsize,
  execute at begin node={\hyphenpenalty=10000},
  tahap/.style={draw, rounded corners, align=center, minimum height=1.0cm,
                text width=3.0cm, inner sep=3pt},
  peran/.style={draw, rounded corners, align=center, minimum height=0.95cm,
                text width=2.4cm, inner sep=3pt, fill=blue!8},
  panah/.style={-{Stealth[length=2mm]}, thick}
]
  % Baris 1: dari tanda pertama sampai dampak dihentikan
  \node[tahap, fill=blue!8] (deteksi) at (0, 0)
    {Deteksi\\{\tiny\mdseries \textit{alert} laju pembakaran SLO}};
  \node[tahap, fill=blue!8] (deklarasi) at (4, 0)
    {Deklarasi dan triase\\{\tiny\mdseries komandan insiden ditunjuk}};
  \node[tahap, fill=orange!12] (mitigasi) at (8, 0)
    {Mitigasi\\{\tiny\mdseries hentikan dampak bagi pasien lebih dahulu}};

  % Baris 2: dari pemulihan sampai pelajaran dijalankan
  \node[tahap, fill=blue!8] (pulih) at (8, -2)
    {Pemulihan\\{\tiny\mdseries layanan kembali memenuhi SLO}};
  \node[tahap, fill=orange!12] (postmortem) at (4, -2)
    {\textit{Postmortem}\\{\tiny\mdseries tanpa menyalahkan, paling lambat 5 hari kerja}};
  \node[tahap, fill=green!10] (lanjut) at (0, -2)
    {Tindak lanjut\\{\tiny\mdseries setiap butir punya pemilik dan tenggat}};

  \draw[panah] (deteksi) -- (deklarasi);
  \draw[panah] (deklarasi) -- (mitigasi);
  \draw[panah] (mitigasi) -- (pulih);
  \draw[panah] (pulih) -- (postmortem);
  \draw[panah] (postmortem) -- (lanjut);
  \draw[panah, dashed] (lanjut) -- node[left, font=\tiny, align=right]
    {mencegah\\terulang} (deteksi);

  % Baris 3: peran selama insiden
  \node[draw, rounded corners, fill=praditagreen!15, minimum height=0.6cm,
        text width=11.0cm, align=center, inner sep=3pt] (judulperan) at (4, -3.75)
    {\textbf{Peran yang dipisahkan selama insiden berlangsung}};
  \node[peran] (komandan) at (-0.3, -5.05)
    {Komandan insiden\\{\tiny\mdseries memimpin dan memutuskan}};
  \node[peran] (operasi) at (2.57, -5.05)
    {Pelaksana operasi\\{\tiny\mdseries satu-satunya yang mengubah sistem}};
  \node[peran] (komunikasi) at (5.43, -5.05)
    {Komunikasi\\{\tiny\mdseries mengabari direksi dan pihak terdampak}};
  \node[peran] (rencana) at (8.3, -5.05)
    {Perencanaan\\{\tiny\mdseries mencatat linimasa dan mengatur sif}};
\end{tikzpicture}
